%!TEX root=../mythesis.tex
% Chapter 6 -- AlphaQD

\chapter{AlphaQD: Quality-Diversity Reinforcement Learning for Formulaic Alpha Mining} % Main chapter title
\chaptermark{AlphaQD for Alpha Mining}
\label{ch:alphaqd}

%----------------------------------------------------------------------------------------
\section{Introduction}
\label{sec:alphaqd_intro}

The two preceding chapters made trading \emph{decisions}: an agent observed the market and chose how to adjust its position. This chapter changes the task.\footnote{This chapter is based on the AlphaQD manuscript, which is currently under review.} Rather than acting in the market, the agent must \emph{generate} the predictive signals from which a portfolio is later built. The signals are \emph{formulaic alphas}: short symbolic expressions whose cross-sectional value at each timestamp is treated as a forecast of future returns~\cite{kakushadze2016101}. Each alpha is a tree of arithmetic, cross-sectional, and time-series operators over a small vocabulary of price and volume features. A portfolio is assembled from a \emph{set} of such alphas.

Formulaic alphas remain widely used in low-frequency quantitative investment because they are readable and can be combined into complementary sets. Automating their discovery is, however, a hard exploration problem under three simultaneous pressures: \emph{quality} (the supervisory signal is the information coefficient of a finished formula on noisy data); \emph{diversity} (a near-duplicate adds little diversification value); and \emph{generalisation} (the training window is bounded and the test market is non-stationary).

\paragraph{Motivation.} Pool-based reinforcement-learning methods such as AlphaGen~\cite{yu2023alphagen} score a candidate by its contribution to an evolving pool. The resulting target changes with the pool, which may complicate optimisation. \fref{fig:alphaqd_reward_drift} illustrates this effect for the correlation-penalised surrogate in \eqref{eq:alphaqd_alphagen}; the theoretical comparison in \sref{sec:alphaqd_theory} concerns this surrogate class rather than the published AlphaGen system.

\begin{figure}[!htb]
\centering
\includegraphics[width=0.8\textwidth]{Chapter6/alphaqd_reward_drift}
\caption[Reward drift under two reward functions.]{Per-candidate reward across pool-insertion steps under two reward functions: the pool-dependent correlation-penalised reward of \eqref{eq:alphaqd_alphagen}, an AlphaGen-style objective (left), and the AlphaQD terminal reward (right). Rows are twelve probe alphas; columns are pool-insertion steps; colour encodes the reward (separate scale per panel). The pool-coupled reward shifts as the pool grows; the AlphaQD reward does not.}
\label{fig:alphaqd_reward_drift}
\end{figure}

\paragraph{Remedy.} AlphaQD removes this coupling in three stages (\fref{fig:alphaqd_pipeline}), following the decompose--specialise--aggregate pattern with the aggregation mechanism fully detached from the reward.
\begin{enumerate}
\item \textbf{GFlowNet exploration.} A relational-graph GFlowNet samples symbolic expressions and is trained by trajectory balance against a terminal reward that depends \emph{only} on the candidate's training-window information ratio, so the learning target is invariant to later archive updates.
\item \textbf{Return-source MAP-Elites archiving.} Each candidate is backtested into a daily IC series; its information ratio scores quality, and the normalised IC series is its behaviour descriptor---an empirical proxy for the formula's return source---placing it in a PCA-discretised MAP-Elites grid~\cite{mouret2015mapelites,pugh2016quality}.
\item \textbf{Top-$K$ cell selection.} One high-quality representative per occupied cell is exported, and the top $K$ cells form the deployed pool.
\end{enumerate}
Against eleven baselines on the reported Chinese A-share and U.S.\ panels, AlphaQD is competitive on the signal metrics. On A-shares it has the highest annualised return among the learned formula-mining methods; among the four generative/RL pool-style methods it has the lowest turnover on both panels.

\begin{figure}[!htb]
\centering
\includegraphics[width=\textwidth]{Chapter6/alphaqd_pipeline}
\caption[The AlphaQD pipeline.]{The three stages of AlphaQD. \emph{Left:} a relational-graph GFlowNet samples a legal symbolic expression and is trained by trajectory balance using a candidate-only reward. \emph{Middle:} the expression is backtested into a daily IC series and inserted into a PCA-discretised MAP-Elites grid. \emph{Right:} the top $K$ occupied cells export their representatives for evaluation with a shared ridge combiner.}
\label{fig:alphaqd_pipeline}
\end{figure}

%----------------------------------------------------------------------------------------
\section{Problem Formulation}
\label{sec:alphaqd_formulation}

Let $T$ be the number of trading days in the training window, $N$ the cross-sectional universe size, and $X\in\mathbb{R}^{T\times N\times d}$ the per-asset feature tensor (open, high, low, close, volume, VWAP, and short-window returns). A formulaic alpha is a function $\alpha:\mathbb{R}^{T\times N\times d}\to\mathbb{R}^{T\times N}$ realised as a reverse-Polish-notation (RPN) token sequence over a typed grammar $\mathcal{G}$, with an action mask rejecting any prefix that cannot complete into a legal formula. Given the $h$-day forward return $r^{(h)}_{t,i}=\mathrm{close}_{t+h,i}/\mathrm{close}_{t,i}-1$ (we use $h=20$), the daily information coefficient and information ratio are, restating the definitions of \sref{sec:bg_metrics} in this chapter's notation~\cite{grinold2000active},
\begin{equation}
\label{eq:alphaqd_ic}
\mathrm{IC}_t(\alpha) = \corr(\alpha_t, r^{(h)}_t),
\qquad
\mathrm{ICIR}(\alpha) = \overline{\mathrm{IC}(\alpha)}\,\big/\,\hat\sigma(\mathrm{IC}(\alpha)).
\end{equation}

\paragraph{Mining goal.} Let $\mathcal{F}$ be the set of legal formulas. The output is a bounded set $\mathcal{P}\subset\mathcal{F}$ whose members, once combined by a fixed downstream combiner, should predict returns on data not seen during mining. The goal is therefore not the per-formula maximiser $\argmax_{\alpha\in\mathcal{F}}\mathrm{ICIR}(\alpha)$, which a single high-ICIR formula would solve, but the joint problem
\begin{equation}
\label{eq:alphaqd_goal}
\argmax_{\mathcal{P}\subset\mathcal{F}}\;\mathbb{E}_{\mathrm{oos}}\!\left[\mathrm{IC}(g(\mathcal{P}))\right]
\quad\text{subject to}\quad |\mathcal{P}|\le K,
\end{equation}
where $g(\cdot)$ is the fixed combiner, $K$ is the pool-size budget, and the expectation is taken over out-of-sample data. The archive-cell map introduced later is denoted separately by $c(\alpha)$. Equation~\eqref{eq:alphaqd_goal} is an idealised target rather than an operational objective: no mining procedure can evaluate the out-of-sample expectation, so search must optimise training-window proxies, and the relation between those proxies and \eqref{eq:alphaqd_goal} is precisely the generalisation question that \sref{sec:alphaqd_theory} addresses. The target imposes two structural requirements: each retained alpha must be of high quality, and members must capture different \emph{return sources} so their combination generalises. We use the centred, normalised daily IC series as the empirical proxy for the return source carried by $\alpha$.

\paragraph{Contrast with AlphaGen.} AlphaGen approximates \eqref{eq:alphaqd_goal} by a boosting relaxation: a pool $\mathrm{Pool}_t$ is grown one element at a time, and the RL reward is the increment in the combined pool's predictive performance after the combination model is refitted with the candidate included. Throughout this chapter we work with a correlation-penalised \emph{variant} that captures the pool coupling of this objective in closed form,
\begin{equation}
\label{eq:alphaqd_alphagen}
R_t(\alpha) = \mathrm{ICIR}(\alpha) - \lambda\max_{\alpha_j\in\mathrm{Pool}_t}\bigl|\corr(\mathrm{IC}(\alpha),\mathrm{IC}(\alpha_j))\bigr|.
\end{equation}
with the correlation penalty defined as zero when $\mathrm{Pool}_t$ is empty. Equation~\eqref{eq:alphaqd_alphagen} is a surrogate for---not a restatement of---the original AlphaGen reward, and the analyses and instrumentation that follow (\fref{fig:alphaqd_reward_drift}, \sref{sec:alphaqd_theory}) address this variant and the pool-coupled reward class it represents rather than the published system. What matters here is the property the variant shares with the original: every pool update changes $R_t(\cdot)$, so the policy is trained against a moving target.

\paragraph{Static formulation.} We retain the per-formula MDP $(\mathcal{S},\mathcal{A},P,R,\gamma=1)$, in which $s_t$ is a partial RPN prefix and $a_t$ the next token under the legal-action mask, but we replace \eqref{eq:alphaqd_alphagen} with a terminal reward that depends only on the candidate and the fixed training window. Diversity is moved out of the reward and into a separate selection rule over a behaviour descriptor of $\alpha$. The reward does not change as the archive grows, removing archive-induced reward drift. The selection rule operates over a bounded number of behaviour cells, yielding a conditional finite-class comparison in \sref{sec:alphaqd_theory}.

%----------------------------------------------------------------------------------------
\section{Method}
\label{sec:alphaqd_method}

The three stages exchange candidates and scalar scores but never a pool-dependent reshape of the reward. Algorithm~\ref{alg:alphaqd} ties them together.

\subsection{GFlowNet Exploration}
\label{sec:alphaqd_generator}

\paragraph{Search surface.} A state $s_t$ is the partial RPN prefix and an action $a_t$ the next token; a legal-action mask, derived from stack depth, type consistency, window-parameter constraints, and a terminator rule, confines the search to syntactically and dimensionally valid expressions. Each prefix is materialised as a typed directed graph---nodes tagged by operator, operand, constant, and window type, edges encoding operator-to-argument relations---and encoded by a two-layer relational graph convolutional network~\cite{schlichtkrull2018rgcn} of hidden width $128$. The pooled node embeddings feed two heads, $\log P_F(\cdot\mid s_t)$ over legal next tokens and $\log P_B(\cdot\mid s_{t+1})$ over legal predecessors, and a scalar $\log Z$.

\paragraph{Terminal reward.} For a completed expression $\alpha$, AlphaQD evaluates its factor values on the training window and cross-sectionally standardises both the factor and the forward-return target to zero mean and unit population standard deviation on each date. Values that remain non-finite after this operation are replaced by zero. The implementation does not winsorise or rank-normalise the factor before computing Pearson IC; this differs from the winsorisation-and-rank-normalisation description in the original method write-up. It then computes $\mathrm{IC}(\alpha)$ and $\mathrm{ICIR}(\alpha)$ via \eqref{eq:alphaqd_ic}, and sets
\begin{equation}
\label{eq:alphaqd_reward}
R(\alpha) = \max\!\bigl\{\exp(\beta_q\,\mathrm{ICIR}(\alpha)),\,\exp(-10)\bigr\},
\end{equation}
with quality temperature $\beta_q$ ($\beta_q=0.5$ in the headline runs; \sref{sec:alphaqd_ablation}) and an exponential floor for an invalid expression, an evaluation failure, or a non-finite or degenerate ICIR. Individual non-finite entries in an otherwise usable daily IC series are zero-filled before the ICIR is formed. The floor also binds for any evaluable candidate with $\mathrm{ICIR}(\alpha)\le -10/\beta_q$, so below that level the generator does not distinguish syntactically valid but extremely poor formulas from unevaluable ones. Equation~\eqref{eq:alphaqd_reward} reads neither the archive nor any pool statistic---not coverage, novelty, cell identity, incumbent quality, or acceptance status---so for a fixed training window and fixed $\beta_q$ the reward of a candidate is invariant to later archive updates.

\paragraph{Trajectory balance.} The generator is trained by trajectory balance~\cite{malkin2022trajectorybalance,bengio2021gflownet}: for a terminal trajectory $\tau$ of token length $L$ (distinct from the trading-day count $T$) with formula $\alpha(\tau)$,
\begin{equation}
\label{eq:alphaqd_tb}
\begin{aligned}
\mathcal{L}_{\mathrm{TB}}(\tau)
=\Biggl(&\log Z+\sum_{t=0}^{L-1}\log P_F(a_t\mid s_t)\\[-0.2em]
&-\sum_{t=0}^{L-1}\log P_B(s_t\mid s_{t+1})
-\log R(\alpha(\tau))\Biggr)^2.
\end{aligned}
\end{equation}
For a finite acyclic construction graph with positive rewards, sufficient model capacity, and an exact trajectory-balance optimum, the terminal sampling probability is proportional to $R(\alpha)$~\cite{bengio2021gflownet}. The trained model only approximates this objective. Failed evaluations receive the floor reward and contribute a finite training signal.

\subsection{Return-Source MAP-Elites Archiving}
\label{sec:alphaqd_archive}

\paragraph{Quality and behaviour descriptor.} The IC series $\mathrm{IC}(\alpha)\in\mathbb{R}^T$ supplies both the quality score $\mathrm{ICIR}(\alpha)$ and a behaviour descriptor that summarises \emph{when} the signal is strong. AlphaQD centres the IC series and normalises it to unit length up to the $\varepsilon$-regularisation in the denominator,
\begin{equation}
\label{eq:alphaqd_shape}
\mathbf{r}(\alpha) = \frac{\mathbf{x}(\alpha)-\overline{\mathbf{x}(\alpha)}\mathbf{1}}{\|\mathbf{x}(\alpha)-\overline{\mathbf{x}(\alpha)}\mathbf{1}\|_2+\varepsilon},
\qquad \mathbf{x}(\alpha)=\mathrm{IC}(\alpha),
\end{equation}
removing level and scale (absorbed into ICIR) and keeping only the relative temporal pattern. By construction $\mathbf{r}(\alpha)$ is invariant to any positive affine transform of the IC series---exactly in the $\varepsilon\to 0$ limit, and up to a negligible $O(\varepsilon)$ term for the machine-precision $\varepsilon$ used in practice (Lemma~\ref{lem:alphaqd_affine})---so it isolates which dates the formula explains, not how strong its edge is; we treat it as the empirical proxy for the formula's return source. The shape is projected by an incremental PCA basis $W$ (canonical dimension $d_b=2$) to $\mathbf{z}(\alpha)=W^\top\mathbf{r}(\alpha)$, then discretised by per-axis quantile bin edges into a $B\times B$ grid ($B=10$, so $100$ cells), giving the cell index $c(\alpha)$. The initial basis and edges are fitted after the first $2{,}000$ \emph{finite candidates}, not after $2{,}000$ episodes. Those candidates are kept pending and are offered to the newly created cells retrospectively after the fit. Thereafter only IC series of \emph{accepted elites} enter a FIFO buffer of at most $1{,}024$ series. In the current AIFML experiment grid, every $K_{\mathrm{pca}}=100$ completed episodes the implementation checks whether at least $32$ accepted series are available; if so, it applies \texttt{IncrementalPCA.partial\_fit} to the buffer, moves each quantile edge towards its new buffer quantile with EMA decay $0.99$, and reprojects every stored elite into the updated cells. When reprojection makes a cell exceed capacity, only the $\kappa$ elites with largest $\mathrm{ICIR}+\mathrm{PFS}$ are retained. Thus both coordinates and membership can change throughout a $2{,}000$-episode run; the descriptor map is not frozen after warm-up.

\paragraph{Per-cell score and update.} For a positive-ICIR candidate, PFS is the mean \emph{ICIR-retention ratio} over the successful perturbations,
\[
\mathrm{PFS}(\alpha)=\operatorname{clip}_{[0,2]}\!\left(\frac{1}{|\mathcal S_\alpha|}\sum_{\widetilde\alpha\in\mathcal S_\alpha}\frac{\mathrm{ICIR}(\widetilde\alpha)}{\mathrm{ICIR}(\alpha)}\right).
\]
The implementation attempts five perturbations. Each rolling-window node is independently replaced by another legal window. It first searches within the configured $\pm10\%$ band, but the legal set is $\{10,20,30,40,50\}$; when that band contains no alternative it samples from the two nearest legal alternatives. Failed or non-finite perturbations are omitted, and PFS is set to $1$ when none succeeds, including formulas with no rolling-window node. Values are clipped to $[0,2]$.

PFS is computed lazily. If a cell is full, a cheap pre-screen first requires the candidate's ICIR to be at least the smallest incumbent ICIR up to $10^{-6}$. A non-positive-ICIR candidate is rejected from the archive in the current configuration and does not update the post-warm-up PCA buffer. A non-full cell accepts an eligible candidate without a dominance check. In a full cell, a candidate that dominates one or more incumbents evicts the dominated incumbent with smallest $\mathrm{ICIR}+\mathrm{PFS}$; a candidate dominated by any incumbent is rejected; otherwise it replaces the incumbent with the smallest sum only when its own sum is strictly larger. This is a bounded score-based cell, not a maintained Pareto front: insertion into a non-full cell can retain dominated members, and one dominance replacement need not remove every dominated member. The quota still prevents a dense region from occupying more than $\kappa$ archive slots per cell.

\subsection{\texorpdfstring{Top-$K$ Cell Selection}{Top-K Cell Selection}}
\label{sec:alphaqd_export}

For each occupied cell, the representative is the elite with the highest training-window ICIR; the export code has no explicit PFS tie-break, so an exact ICIR tie retains iteration order. Let $K^\star=\min(10,5\kappa)$, i.e.\ $5\kappa$ capped at $10$: $K^\star=10$ under the canonical $\kappa=3$, matching the pool-style baselines, but $K^\star=5$ under $\kappa=1$, so the rule couples the exported pool size to the per-cell capacity---a coupling that matters when interpreting the $\kappa$ axis of the sensitivity grid in \sref{sec:alphaqd_ablation}. AlphaQD ranks occupied cells by representative ICIR and exports the top $K^\star$ as the pool $\mathcal{P}^\star$. The main comparison then uses the shared closed-form ridge combiner~\cite{hoerl1970ridge},
\[
\widehat{\mathbf w}=(X^\top X+I)^{-1}X^\top\mathbf y,
\]
fitted to tradable validation-split asset--date rows with finite targets. There is no intercept or feature standardisation; missing factor values are replaced by zero. The regularisation coefficient is fixed at $\lambda=1.0$, and the resulting weights are held fixed on the test set. Positive-ICIR-proportional weights are written with the native AlphaQD artifact, but they are not the weights used by the unified-ridge headline table. Because the per-cell representative inherits the positive-ICIR admission rule and the cells partition formulas by behaviour, the export is designed to encourage positive standalone training-window edge and lower redundancy; whether the exported pool meets the formal premises of \sref{sec:alphaqd_theory} out of sample is not guaranteed (see the discussion after Theorem~\ref{thm:alphaqd_ensemble}).

\begin{algorithm}[!htb]
\caption{AlphaQD: three-stage training and export}
\label{alg:alphaqd}
\begin{algorithmic}[1]
\State Initialise the RGCN encoder, $\log P_F$, $\log P_B$, $\log Z$, an empty MAP-Elites archive, and an empty PCA buffer.
\For{episode $= 1,\dots,E$}
  \State \textbf{Stage 1:} sample a batch of terminal trajectories under $P_F$ and the legal-action mask.
  \For{each terminal trajectory $\tau$ with formula $\alpha$}
    \State Evaluate $\alpha$ on the training window; compute $\mathrm{IC}(\alpha)$, $\mathrm{ICIR}(\alpha)$, and $R(\alpha)$ via \eqref{eq:alphaqd_reward}.
    \State \textbf{Stage 2:} update the descriptor and bounded archive as follows.
    \If{the first $2{,}000$ finite candidates have not yet been collected}
      \State Append $\mathbf{r}(\alpha)$ to the warm-up set and retain $\alpha$ as pending; at the threshold, fit $W$ and the initial quantile edges and offer all pending positive-ICIR candidates to their cells.
    \ElsIf{$\mathrm{ICIR}(\alpha)>0$ and the cell ICIR pre-screen passes}
      \State Compute PFS from successful legal-window perturbations and apply the bounded cell update rule.
      \If{$\alpha$ is accepted}
        \State Append its IC series to the $1{,}024$-element FIFO PCA buffer.
      \EndIf
    \EndIf
    \State Accumulate $\mathcal{L}_{\mathrm{TB}}(\tau)$ from \eqref{eq:alphaqd_tb}.
  \EndFor
  \State Update $\{\log P_F,\log P_B,\log Z\}$ on the accumulated loss.
  \If{the episode count is a multiple of $K_{\mathrm{pca}}=100$ and the FIFO buffer has at least $32$ series}
    \State Incrementally update $W$, move the quantile edges by EMA, reproject all elites, and trim over-capacity cells by $\mathrm{ICIR}+\mathrm{PFS}$.
  \EndIf
\EndFor
\State \textbf{Stage 3:} take the highest-ICIR elite per cell as its representative; rank cells by representative ICIR; export the top $K^\star$ as $\mathcal{P}^\star$ and fit the shared $\lambda=1$ ridge combiner on validation rows.
\end{algorithmic}
\end{algorithm}

%----------------------------------------------------------------------------------------
\section{Conditional Analytical Properties}
\label{sec:alphaqd_theory}

This section gives conditional analytical results, with proofs inline. The deployed descriptor map is data-dependent and changes at the periodic PCA updates described in \sref{sec:alphaqd_archive}; none of the results below treats that moving partition as fixed unless the premise says so explicitly. One piece of notation is needed by Theorem~\ref{thm:alphaqd_stationary}: for a cell index $c=(b_1,\dots,b_{d_b})$ with $b_j\in\{0,\dots,B-1\}$, the \emph{anti-cell} is $\mathrm{anti}(c)=(B-1-b_1,\dots,B-1-b_{d_b})$, the \emph{orbit} of $c$ is the unordered pair $\mathcal{O}(c)=\{c,\mathrm{anti}(c)\}$, and, by a mild abuse of notation, $\mathcal{O}$ also denotes the set of all orbits. A formula and its sign-flip have oppositely signed IC shapes, but empirical quantile edges---especially after EMA updates---need not be symmetric, so antipodal cell membership is an idealised condition rather than a deployed invariant. Throughout, the discussion after each theorem distinguishes such fixed-map statements from the implementation.

\subsection{Decoupling Quality from Shape}
\label{sec:alphaqd_decouple}

The first result formalises why decoupling quality from the shape descriptor is sound.

\begin{lemma}[Affine invariance of the shape descriptor]
\label{lem:alphaqd_affine}
For a candidate $\alpha$, write $\mathbf{v}=\mathbf{x}(\alpha)-\overline{\mathbf{x}(\alpha)}\mathbf{1}$ and suppose $\|\mathbf{v}\|_2>0$. For any positive affine transform $\mathbf{y}=a\,\mathbf{x}(\alpha)+b\mathbf{1}$ with $a>0$, let $\mathbf{r}_{\mathbf{y}}$ denote the shape vector \eqref{eq:alphaqd_shape} built from $\mathbf{y}$. Then $\mathbf{r}_{\mathbf{y}}=\mathbf{r}(\alpha)$ exactly in the limit $\varepsilon\to 0$, and for $\varepsilon>0$,
\[
\bigl\|\mathbf{r}_{\mathbf{y}}-\mathbf{r}(\alpha)\bigr\|_2
=\frac{\varepsilon\,|a-1|\,\|\mathbf{v}\|_2}{(a\|\mathbf{v}\|_2+\varepsilon)(\|\mathbf{v}\|_2+\varepsilon)}
=O\!\bigl(\varepsilon/\|\mathbf{v}\|_2\bigr).
\]
\end{lemma}
\begin{proof}
The mean satisfies $\overline{\mathbf{y}}=a\,\overline{\mathbf{x}(\alpha)}+b$, so $\mathbf{y}-\overline{\mathbf{y}}\mathbf{1}=a\mathbf{v}$ and $\mathbf{r}_{\mathbf{y}}=a\mathbf{v}/(a\|\mathbf{v}\|_2+\varepsilon)$. Hence
\[
\mathbf{r}_{\mathbf{y}}-\mathbf{r}(\alpha)
=\Bigl(\frac{a}{a\|\mathbf{v}\|_2+\varepsilon}-\frac{1}{\|\mathbf{v}\|_2+\varepsilon}\Bigr)\mathbf{v}
=\frac{\varepsilon\,(a-1)}{(a\|\mathbf{v}\|_2+\varepsilon)(\|\mathbf{v}\|_2+\varepsilon)}\,\mathbf{v},
\]
whose norm is the stated quantity. It vanishes as $\varepsilon\to 0$ (and for $a=1$); with $\varepsilon=0$ the normalisation cancels $a>0$ exactly.
\end{proof}

Thus, up to an error of order $\varepsilon$---which is fixed at machine-precision scale while $\|\mathbf{v}\|_2$ is bounded away from zero for any non-constant IC series---the descriptor encodes only which dates were strong; the mean and scale are absorbed into the quality score. For the same reason the descriptor is only approximately of unit norm, $\|\mathbf{r}(\alpha)\|_2=\|\mathbf{v}\|_2/(\|\mathbf{v}\|_2+\varepsilon)<1$, and ``unit-normalised'' is used in this $O(\varepsilon)$ sense throughout. Two alphas with the same temporal pattern but different overall edge occupy the same cell rather than leaking into different ones; a zero-variance IC series ($\|\mathbf{v}\|_2=0$) has no meaningful shape and is handled by the reward floor of \eqref{eq:alphaqd_reward} rather than by the descriptor.

\subsection{Reward Stationarity}
\label{sec:alphaqd_stationary}

The second result is the structural reason the static reward survives even archive-aware exploration extensions.

\begin{theorem}[Stationary conditional reward under a fixed descriptor map]
\label{thm:alphaqd_stationary}
Let the per-alpha reward be $R_q(\alpha)=\exp(\beta_q\,\mathrm{ICIR}(\alpha))$, and assume for this conditional extension that the PCA basis $W$ and the bin edges entering $c(\cdot)$ are fixed. For any fixed target orbit $\mathcal{O}^\star$ and any kernel $k:\mathcal{O}\times\mathcal{O}\to\mathbb{R}_{\ge 0}$ independent of the archive state and training step, the conditional reward $R_{\mathcal{O}^\star}(\alpha)=R_q(\alpha)\,k(\mathcal{O}(c(\alpha)),\mathcal{O}^\star)$ is a stationary terminal reward.
\end{theorem}
\begin{proof}
For each fixed $\mathcal{O}^\star$, $R_{\mathcal{O}^\star}(\cdot)$ depends on $\alpha$ only through $\mathrm{ICIR}(\alpha)$ on the fixed training window and through the fixed PCA basis and bin edges entering $c(\cdot)$. The kernel $k$ is fixed and $\mathcal{O}^\star$ is a parameter, not a function of training-time state. The reward therefore has no time-dependent component. A scheduler that varies $\mathcal{O}^\star$ across episodes changes only the per-episode prior over conditions, never the conditional reward itself.
\end{proof}

The fixed-map premise is \emph{not} supplied by the deployed schedule: the current code updates the PCA basis and edges every $100$ episodes when its buffer is large enough. The implementation uses the unconditional reward ($k\equiv 1$), whose stationarity is independent of the descriptor map and therefore survives those updates. The conditional result applies only to a hypothetical descriptor-biased extension with a fixed map; it is not a theorem about the deployed rolling-PCA partition.

\subsection{Search Complexity and a Boosting Counterexample}
\label{sec:alphaqd_complexity}

The next two results contrast AlphaQD's quality-diversity ensemble with boosting-style residual fitting in the abstract form used by the pool-coupled rewards of \sref{sec:alphaqd_formulation}; they concern this class of procedures, not the specific published AlphaGen system.

\begin{theorem}[Residual-edge noise floor]
\label{thm:alphaqd_boost}
Let $\mathcal{F}=\{\alpha_1,\dots,\alpha_M\}$ be a finite candidate set. After $t-1$ boosting rounds with residual $r_{t-1}$, define the population residual edge and its training-window estimate by
\[
\gamma_t(\alpha)=\mathbb{E}[\langle\alpha,r_{t-1}\rangle],
\qquad
\hat\gamma_t(\alpha)=\gamma_t(\alpha)+\xi_{t,\alpha}.
\]
Assume that each deviation is sub-Gaussian: for a scale parameter $\sigma$ and effective sample size $n_{\mathrm{eff}}$,
\[
\mathbb{P}(|\xi_{t,\alpha}|>u)
\le 2\exp\!\left(-\frac{n_{\mathrm{eff}}u^2}{2\sigma^2}\right).
\]
Let $\alpha_t=\argmax_{\alpha\in\mathcal{F}}\hat\gamma_t(\alpha)$, $\Delta_t=\sup_\alpha\gamma_t(\alpha)$, and $\eta_M=\sigma\sqrt{2\log(2M/\delta)/n_{\mathrm{eff}}}$. Then, with probability at least $1-\delta$, $\gamma_t(\alpha_t)\ge\Delta_t-2\eta_M$. When $\Delta_t\le 2\eta_M$, the bound is non-positive.
\end{theorem}
\begin{proof}
Uniform sub-Gaussian concentration on the finite class $\mathcal{F}$ gives $\sup_\alpha|\hat\gamma_t(\alpha)-\gamma_t(\alpha)|\le\eta_M$ with probability at least $1-\delta$. Writing $\alpha_t^\star=\argmax_\alpha\gamma_t(\alpha)$, $\gamma_t(\alpha_t)\ge\hat\gamma_t(\alpha_t)-\eta_M\ge\hat\gamma_t(\alpha_t^\star)-\eta_M\ge\gamma_t(\alpha_t^\star)-2\eta_M=\Delta_t-2\eta_M$, where the middle step uses the empirical optimality of $\alpha_t$.
\end{proof}

With $|\mathcal{F}|>10^{20}$ and $n_{\mathrm{eff}}$ bounded by the training-window days, $\eta_M$ is large relative to plausible late-round residual edges, so the guarantee $\Delta_t-2\eta_M$ turns non-positive and hence vacuous: the empirical residual winner retains no assurance of positive population edge. This is a statement about the emptiness of the guarantee, not an upper bound on the winner's edge; establishing that late-round winners are in fact noise-dominated would require a matching upper bound or an explicit winner's-curse argument, which we do not attempt.

\begin{theorem}[Search-complexity contrast]
\label{thm:alphaqd_complexity}
A $K$-step boosting procedure over $\mathcal{F}$ with $|\mathcal{F}|=M$ (here $K$ plays the same retained-set-size role as in \eqref{eq:alphaqd_goal}) has an ordered-sequence hypothesis class of size at most $M^K$, with finite-class complexity $\mathcal{C}_{\mathrm{boost},K}=\sqrt{K\log M/n_{\mathrm{eff}}}$. A quality-diversity rule mapping candidates to at most $G$ occupied cells and exporting one representative per cell has $K$-subset class of size at most $\binom{G}{K}$, with complexity $\mathcal{C}_{\mathrm{QD},K}=\sqrt{K\log(eG/K)/n_{\mathrm{eff}}}$. For $G\ll M$, $\mathcal{C}_{\mathrm{QD},K}\ll\mathcal{C}_{\mathrm{boost},K}$.
\end{theorem}
\begin{proof}
The finite-class generalisation bound scales as $\sqrt{\log|\mathcal{H}|/n_{\mathrm{eff}}}$. Substituting $|\mathcal{H}_{\mathrm{boost},K}|\le M^K$ gives $\log|\mathcal{H}_{\mathrm{boost},K}|\le K\log M$. Substituting $|\mathcal{H}_{\mathrm{QD},K}|\le\binom{G}{K}$ and $\log\binom{G}{K}\le K\log(eG/K)$ gives the QD scaling.
\end{proof}

In the reported grid $G\le 100$ while $|\mathcal{F}|>10^{20}$: the class sizes differ by many orders of magnitude, though the complexity terms themselves differ only by a modest factor (about $3.7$ under $M=10^{20}$, $G=100$, $K=10$). Two caveats govern the reading. First, the comparison is conditional on the realised archive: the count $\binom{G}{K}$ treats the per-cell representatives, the PCA basis, and the bin edges as fixed, whereas all three are data-dependent products of a search that itself screened a number of candidates comparable to what boosting sees, and the ridge weights add a continuous layer on top. Theorem~\ref{thm:alphaqd_complexity} therefore bounds the complexity of the final selection stage given the archive, not the effective complexity of the full pipeline. Second, the two class-size counts are upper bounds on differently structured procedures, so the contrast is indicative rather than a like-for-like risk comparison.

\subsection{Per-Cell Selection Bias}
\label{sec:alphaqd_selection}

When a cell aggregates several candidates and the elite is selected by maximum empirical ICIR, the chosen elite's empirical ICIR is upward-biased relative to its population value---the formula-mining analogue of a multiple-testing winner's curse~\cite{harvey2016cross}. The bounded per-cell quota controls the final, within-cell component of this bias.

\begin{theorem}[Per-cell selection bias under a bounded quota]
\label{thm:alphaqd_selection}
Fix an occupied cell holding at most $\kappa$ admitted elites at selection time, each with population quality $\vartheta_i$ and training-window estimate $\hat\vartheta_i=\vartheta_i+\varepsilon_i$, where each $\varepsilon_i$ is sub-Gaussian with $\mathbb{P}(\varepsilon_i>u)\le\exp(-n_{\mathrm{eff}}u^2/(2\sigma^2))$ for a scale parameter $\sigma$ and effective sample size $n_{\mathrm{eff}}$. With probability at least $1-\delta$, the selection bias of the per-cell representative on raw ICIR satisfies $\max_{i\le\kappa}(\hat\vartheta_i-\vartheta_i)\le\sigma\sqrt{2\log(\kappa/\delta)/n_{\mathrm{eff}}}$.
\end{theorem}
\begin{proof}
A union bound on the $\kappa$ sub-Gaussian tails gives $\mathbb{P}(\max_{i\le\kappa}\varepsilon_i>u)\le\kappa\exp(-n_{\mathrm{eff}}u^2/(2\sigma^2))$; setting the right-hand side to $\delta$ yields $u=\sigma\sqrt{2\log(\kappa/\delta)/n_{\mathrm{eff}}}$.
\end{proof}

Under the canonical $\kappa=3$, $\delta=0.05$ the bound is $\approx 2.86\,\sigma/\sqrt{n_{\mathrm{eff}}}$. Two caveats bound its reach. First, the deployed ranking uses raw ICIR with PFS only as a tie-break (\sref{sec:alphaqd_archive}); the five PFS perturbations share the same training window and are neither independent nor same-mean estimates of the elite's quality, so no $\sqrt{1+m}$-style tightening of the bound is claimed here---any such tightening would be heuristic. Second, and more importantly, the theorem conditions on the $\kappa$ stored elites as a fixed set, whereas they are themselves the survivors of an empirical-score tournament over every candidate the cell ever received. The quota therefore bounds the bias of the final within-cell $\argmax$, but it does not remove the selection bias accumulated during the search---the search-stage multiple-testing winner's curse~\cite{harvey2016cross} lies outside the theorem's scope.

\subsection{Ensemble ICIR Lower Bound}
\label{sec:alphaqd_ensemble}

The final result is an abstract lemma about equal-weight ensembles; its relation to the deployed pool is discussed below.

\begin{theorem}[Ensemble ICIR lower bound]
\label{thm:alphaqd_ensemble}
Let $X_1,\dots,X_K$ be daily IC variables satisfying $\mathbb{E}[X_i]\ge\mu_{\min}>0$, $\Var(X_i)\le\sigma^2$, and $\corr(X_i,X_j)\le\rho_{\max}$ for all $i\ne j$, for some $\rho_{\max}\in[0,1]$. Then the equal-weight ensemble $X_{\mathrm{ens}}=\frac{1}{K}\sum_i X_i$ obeys
\begin{equation}
\label{eq:alphaqd_ensemble}
\mathrm{ICIR}(X_{\mathrm{ens}}) \ge \frac{\mu_{\min}\sqrt{K}}{\sigma\sqrt{1+(K-1)\rho_{\max}}}.
\end{equation}
\end{theorem}
\begin{proof}
By linearity, $\mathbb{E}[X_{\mathrm{ens}}]\ge\mu_{\min}$. For the variance,
\begin{align*}
\Var(X_{\mathrm{ens}})
&=\frac{1}{K^2}\sum_i\Var(X_i)
  +\frac{1}{K^2}\sum_{i\ne j}\Cov(X_i,X_j)\\
&\le\frac{\sigma^2}{K^2}\bigl[K+K(K-1)\rho_{\max}\bigr]\\
&=\frac{\sigma^2\bigl(1+(K-1)\rho_{\max}\bigr)}{K}.
\end{align*}
so $\mathrm{ICIR}(X_{\mathrm{ens}})=\mathbb{E}[X_{\mathrm{ens}}]/\sqrt{\Var(X_{\mathrm{ens}})}\ge\mu_{\min}\sqrt{K}/(\sigma\sqrt{1+(K-1)\rho_{\max}})$.
\end{proof}

The proof also uses $\rho_{\max}\ge 0$: with a negative correlation cap the covariance inequality would reverse. Theorem~\ref{thm:alphaqd_ensemble} is a statement about any family of random variables satisfying its premises; three gaps separate it from a guarantee for the deployed pool. First, the admission rule enforces a positive \emph{training-window} ICIR, which makes the premise $\mathbb{E}[X_i]\ge\mu_{\min}>0$ plausible for the exported elites but does not imply it out of sample. Second, the premise requires a uniform cap $\rho_{\max}$ on all pairwise correlations of the exported components. Spreading elites across PCA cells encourages lower redundancy but does not enforce such a cap, and the archive figure reports no pairwise-correlation statistic. Third, the bound concerns the equal-weight average of the components' daily IC variables, whereas the deployed signal is a ridge-weighted combination of factor values whose IC is not a linear average of the component ICs. Theorem~\ref{thm:alphaqd_boost} shows that boosting-style residual fitting need not preserve even the first premise, so the analogous bound is not available for it by this route; for AlphaQD too, however, the result is motivation for the export design rather than a finite-sample guarantee for the deployed pool.

%----------------------------------------------------------------------------------------
\section{Experimental Setup}
\label{sec:alphaqd_setup}

\paragraph{Datasets.} We evaluate on two markets so the comparison is not anchored to one regime. The implementation is maintained in the AIFML project repository.\footnote{\url{https://github.com/qinmoelei/AIFML}} The Chinese A-share panel (the canonical alpha-mining setting) is the primary benchmark: the experiment record reports BaoStock daily data for $7{,}207$ instruments over $3{,}963$ trading days from 2010-01-04 to 2026-04-30, with front-adjusted OHLCV and $\mathrm{VWAP}=\mathrm{amount}/\mathrm{volume}$. The U.S.\ panel is an independent check with different microstructure and base rates: the record reports a Tiingo Russell-3000-style panel of $2{,}571$ instruments over $4{,}109$ trading days, using Tiingo's adjusted price fields. Tradability masks apply at execution: on A-shares an asset must have finite prices, positive volume, normal trading status, non-ST status, and an open-to-previous-close gap within $\pm 9.7\%$ (approximating price-limit constraints); on the U.S.\ panel, finite prices, positive volume, and a minimum close of one dollar (excluding penny names).

The datasets and preprocessing follow the experimental setting established by AlphaGen~\cite{yu2023alphagen}. Every method within a market receives the same instruments, dates, labels, execution rule, transaction costs, and portfolio construction protocol, enabling a controlled comparison of the alpha-mining methods.

\paragraph{Split and usage.} All methods share the same three-way chronological split: 2010--2021 for training, 2022 for validation, and 2023-01-01 to 2026-04-30 for testing, with forward-return horizon $h=20$ trading days. Following the AlphaGen protocol, the training set is used to search for or learn formulaic alphas. Once a method exports its alpha pool, the validation set is used to fit the shared ridge combiner that maps the pool to a scalar signal. The fitted linear weights are then fixed, and the test set is used only for final performance evaluation. This separation gives every pool-style method the same downstream combiner and isolates the quality of the generated alpha pool from method-specific weighting choices.

\paragraph{Random seeds.} Stochastic methods are evaluated using seeds 0, 1, 2, and 12345. We report the mean and population standard deviation across these runs; deterministic methods have zero run variance.

\paragraph{Compared methods.} We compare against eleven baselines in four families: formulaic (the fixed Alpha101 catalogue~\cite{kakushadze2016101}); supervised (MLP~\cite{murtagh1991mlp}, XGBoost~\cite{chen2016xgboost}, LightGBM~\cite{ke2017lightgbm}); genetic programming (GP single-predictor; GP with mutual-IC filter; GP without filter~\cite{stephens2015gplearn}); and generative/RL (AlphaForge~\cite{shi2025alphaforge}, AlphaGen-PPO~\cite{yu2023alphagen}, AlphaQCM~\cite{zhu2025alphaqcm}, AlphaSAGE~\cite{chen2025alphasage}). Pool-style baselines export ten alphas each, matching AlphaQD's $K^\star=10$.

\paragraph{Shared protocol.} Every pool-style method, including AlphaQD, exports a fixed-size pool. That pool is fed to the same closed-form $\lambda=1$ ridge combiner described in \sref{sec:alphaqd_export}, with weights fitted on validation data and held fixed on the test set. We report signal-layer IC, ICIR, and RankIC~\cite{grinold2000active} together with portfolio-layer annualised net Sharpe~\cite{sharpe1994sharpe}, annualised net return (the arithmetic mean daily net return scaled by $252$), maximum drawdown~\cite{magdon2004maximum} (reported as a positive magnitude; smaller is better), and annualised turnover. The portfolio is long-only top-$50$ rebalanced every twenty trading days, executed at $\mathrm{VWAP}_{t+1}$ on A-shares and $\mathrm{Open}_{t+1}$ on the U.S.\ panel, with round-trip costs of $12.5$ and $3.5$ basis points respectively. Applying this common protocol to every method makes the within-panel ranks the primary comparison.

%----------------------------------------------------------------------------------------
\section{Results and Analysis}
\label{sec:alphaqd_results}

\subsection{Main Results}
\label{sec:alphaqd_main}

\tref{tab:alphaqd_main} reports both panels under the unified-ridge protocol, aggregated over seeds $\{0,1,2,12345\}$ where stochastic.

\begin{table}[!htb]
\centering
\caption[AlphaQD main results.]{Main results on the 2023-01-01 to 2026-04-30 test window under the unified-ridge protocol. Entries are mean $\pm$ population standard deviation over seeds $\{0,1,2,12345\}$ where stochastic; deterministic methods carry zero seed variance. Portfolio metrics are net of round-trip costs on a long-only top-$50$ portfolio rebalanced every twenty trading days; maximum drawdown is reported as a positive magnitude (smaller is better). Per column within each panel, the \fst{best}, \snd{second-best}, and \trd{third-best} entries are marked in bold, underline, and dotted underline respectively (colour retained for screen reading).}
\label{tab:alphaqd_main}
\renewcommand{\arraystretch}{1.15}
\footnotesize
\resizebox{\textwidth}{!}{
\begin{tabular}{llccccccc}
\toprule
Family & Method & IC$\uparrow$ & ICIR$\uparrow$ & RankIC$\uparrow$ & Sharpe$\uparrow$ & AnnRet(\%)$\uparrow$ & MDD(\%)$\downarrow$ & TO/yr$\downarrow$ \\
\midrule
\multicolumn{9}{c}{\textit{Panel A: Chinese A-share}} \\
\midrule
Formulaic & Alpha101 & $0.0487\pm0.0000$ & $0.408\pm0.000$ & $0.082\pm0.000$ & \fst{$0.983\pm0.000$} & \snd{$19.9\pm0.0$} & \fst{$26.1\pm0.0$} & $23.9$ \\
\multirow{3}{*}{Supervised} & MLP & $0.0079\pm0.0017$ & $0.111\pm0.018$ & $0.017\pm0.003$ & $0.319\pm0.038$ & $9.1\pm1.3$ & $49.7\pm1.0$ & $12.9$ \\
 & LightGBM & $0.0189\pm0.0002$ & $0.262\pm0.006$ & $0.035\pm0.000$ & $0.309\pm0.022$ & $10.0\pm0.7$ & $52.7\pm1.4$ & \snd{$12.0$} \\
 & XGBoost & $0.0199\pm0.0000$ & $0.266\pm0.000$ & $0.036\pm0.000$ & $0.236\pm0.000$ & $6.9\pm0.0$ & $50.6\pm0.0$ & \trd{$12.4$} \\
\multirow{3}{*}{Genetic Prog.} & GP (predictor) & $-0.0101\pm0.0132$ & $-0.082\pm0.118$ & $-0.016\pm0.024$ & $0.292\pm0.309$ & $6.4\pm6.1$ & $40.4\pm16.4$ & \fst{$3.8$} \\
 & GP w/ filter & \trd{$0.0652\pm0.0080$} & $0.507\pm0.070$ & \snd{$0.105\pm0.010$} & $0.663\pm0.082$ & \trd{$18.6\pm3.9$} & $39.1\pm5.2$ & $15.1$ \\
 & GP w/o filter & \fst{$0.0703\pm0.0157$} & \fst{$0.581\pm0.064$} & \trd{$0.102\pm0.025$} & $0.501\pm0.573$ & $10.8\pm12.9$ & $35.7\pm9.4$ & $20.2$ \\
\multirow{4}{*}{Generative / RL} & AlphaForge & $0.0636\pm0.0099$ & \trd{$0.553\pm0.030$} & $0.086\pm0.016$ & $0.534\pm0.493$ & $14.4\pm13.4$ & $37.0\pm15.7$ & $19.5$ \\
 & AlphaGen-PPO & $0.0435\pm0.0215$ & $0.495\pm0.219$ & $0.049\pm0.043$ & $0.275\pm0.551$ & $4.1\pm11.3$ & $43.3\pm15.4$ & $22.2$ \\
 & AlphaQCM (FQF) & $0.0206\pm0.0187$ & $0.255\pm0.291$ & $0.020\pm0.035$ & $0.053\pm0.451$ & $1.7\pm12.2$ & $46.8\pm12.2$ & $19.3$ \\
 & AlphaSAGE (GNN) & $0.0552\pm0.0082$ & \snd{$0.555\pm0.142$} & $0.095\pm0.008$ & \trd{$0.730\pm0.050$} & $14.9\pm1.0$ & \snd{$27.4\pm3.9$} & $18.3$ \\
Ours & \textbf{AlphaQD} & \snd{$0.0673\pm0.0093$} & $0.522\pm0.078$ & \fst{$0.112\pm0.011$} & \snd{$0.911\pm0.185$} & \fst{$24.7\pm6.2$} & \trd{$31.5\pm1.7$} & $12.5$ \\
\midrule
\multicolumn{9}{c}{\textit{Panel B: U.S.\ Russell-3000-style universe}} \\
\midrule
Formulaic & Alpha101 & $0.0188\pm0.0000$ & $0.358\pm0.000$ & \snd{$-0.002\pm0.000$} & $1.511\pm0.000$ & $38.2\pm0.0$ & \snd{$15.9\pm0.0$} & $23.0$ \\
\multirow{3}{*}{Supervised} & MLP & $0.0446\pm0.0102$ & \fst{$0.604\pm0.147$} & \trd{$-0.002\pm0.007$} & $1.796\pm0.055$ & $99.8\pm5.0$ & \trd{$16.3\pm0.9$} & $9.1$ \\
 & LightGBM & $0.0404\pm0.0023$ & $0.454\pm0.025$ & $-0.009\pm0.001$ & $2.090\pm0.035$ & $119.1\pm2.3$ & $21.6\pm1.4$ & $8.7$ \\
 & XGBoost & $0.0335\pm0.0000$ & $0.395\pm0.000$ & $-0.006\pm0.000$ & $2.120\pm0.000$ & $129.2\pm0.0$ & $23.4\pm0.0$ & $9.6$ \\
\multirow{3}{*}{Genetic Prog.} & GP (predictor) & $-0.0104\pm0.0458$ & $-0.073\pm0.417$ & \fst{$0.004\pm0.010$} & $1.202\pm0.497$ & $39.4\pm45.8$ & \fst{$15.2\pm5.2$} & \fst{$5.7$} \\
 & GP w/ filter & \snd{$0.0691\pm0.0014$} & $0.512\pm0.069$ & $-0.005\pm0.003$ & \fst{$2.180\pm0.043$} & \fst{$141.6\pm6.4$} & $26.7\pm2.4$ & $9.5$ \\
 & GP w/o filter & \fst{$0.0717\pm0.0031$} & \snd{$0.568\pm0.013$} & $-0.013\pm0.001$ & \trd{$2.164\pm0.047$} & \snd{$141.3\pm5.5$} & $24.7\pm0.8$ & $10.2$ \\
\multirow{4}{*}{Generative / RL} & AlphaForge & $0.0481\pm0.0227$ & $0.430\pm0.181$ & $-0.008\pm0.002$ & $1.856\pm0.525$ & $107.0\pm41.3$ & $23.4\pm1.7$ & $11.6$ \\
 & AlphaGen-PPO & $0.0542\pm0.0082$ & $0.437\pm0.062$ & $-0.010\pm0.002$ & $2.056\pm0.223$ & $129.8\pm15.6$ & $23.3\pm1.1$ & $13.8$ \\
 & AlphaQCM (FQF) & $0.0486\pm0.0295$ & $0.453\pm0.286$ & $-0.002\pm0.009$ & $1.942\pm0.372$ & $103.8\pm46.4$ & $21.1\pm3.4$ & \trd{$7.9$} \\
 & AlphaSAGE (GNN) & $0.0614\pm0.0077$ & $0.539\pm0.068$ & $-0.007\pm0.003$ & $2.111\pm0.094$ & $131.3\pm7.4$ & $24.2\pm1.5$ & $10.5$ \\
Ours & \textbf{AlphaQD} & \trd{$0.0655\pm0.0022$} & \trd{$0.551\pm0.036$} & $-0.010\pm0.001$ & \snd{$2.177\pm0.020$} & \trd{$133.3\pm2.0$} & $22.5\pm0.2$ & \snd{$6.7$} \\
\bottomrule
\end{tabular}}
\end{table}

\paragraph{A-share findings.} On the A-share panel, AlphaQD attains the highest annualised return and the highest RankIC of any method. It is second on IC (behind GP-without-filter) and second on Sharpe (behind only the fixed Alpha101 catalogue). The leading signal-quality group comprises GP-without-filter, AlphaForge, AlphaSAGE, and AlphaQD. Within this group the portfolio columns differentiate more clearly: GP-without-filter and AlphaForge have substantially larger seed-to-seed variation in Sharpe and annualised return, while AlphaSAGE is stable but has an annualised return roughly ten percentage points below AlphaQD and turnover $46\%$ higher. Among the four generative/RL pool-style methods, AlphaQD has the lowest turnover and the smallest standard deviation in maximum drawdown. \fref{fig:alphaqd_nav} plots the cumulative net NAV of the five learned pool-style methods.

\begin{figure}[!htb]
\centering
\includegraphics[width=0.78\textwidth]{Chapter6/alphaqd_portfolio_nav}
\caption[Cumulative NAV on the A-share test window.]{Cumulative net NAV (start $=1.0$) on the A-share test window 2023-01-01 to 2026-04-30 for the five learned pool-style methods. Solid lines are the mean NAV across seeds $\{0,1,2,12345\}$ at each rebalance date; shaded bands span $\pm 1$ population standard deviation; the dashed line marks NAV $=1$. The portfolio is the long-only top-$50$ rebalanced every twenty trading days at $\mathrm{VWAP}_{t+1}$ with $12.5$ bps round-trip cost.}
\label{fig:alphaqd_nav}
\end{figure}

\paragraph{U.S.\ findings.} Under the shared U.S.\ protocol, AlphaQD ranks third on IC, ICIR, and annualised return, and second on net Sharpe and turnover. Among the four generative/RL pool-style methods it also has the smallest reported run-level deviation on every IC-family metric. These relative positions, rather than the absolute level of any one portfolio metric, are the intended cross-method comparison on this panel.

\subsection{Archive Diagnostics}
\label{sec:alphaqd_diagnostics}

The archive snapshot in \fref{fig:alphaqd_archive} contains $87$ occupied cells and $224$ stored elites. The largest representative ICIR annotated in the grid is $0.66$. These figure-derived values are used consistently as the descriptive archive result; the grid visualises broad coverage of the two-dimensional behaviour space, while it does not by itself establish a uniform correlation bound for the exported pool.

\begin{figure}[!htb]
\centering
\includegraphics[width=0.8\textwidth]{Chapter6/alphaqd_archive_heatmap}
\caption[The AlphaQD MAP-Elites archive.]{AlphaQD archive snapshot on the $10\times 10$ behaviour-descriptor grid (vertical axis PC1 bin, horizontal axis PC2 bin). \emph{(a)} number of stored elites per cell; \emph{(b)} representative ICIR per cell, with ten cells outlined. The plotted annotations contain $87$ occupied cells and $224$ elites; the largest annotated representative ICIR is $0.66$.}
\label{fig:alphaqd_archive}
\end{figure}

\subsection{Hyperparameter Sensitivity}
\label{sec:alphaqd_ablation}

A grid over behaviour-descriptor dimension $d_b\in\{1,2,3\}$, per-cell capacity $\kappa\in\{1,3,5\}$, and reward temperature $\beta_q\in\{0.5,1,2,5\}$ probes the sensitivity of the canonical configuration on the validation set. Marginal ICIR declines gently from $d_b=1$ to $d_b=3$, consistent with a higher-dimensional descriptor spreading candidates over more cells; the canonical $d_b=2$ preserves an interpretable two-axis archive while remaining in the favourable range. Capacity $\kappa\ge 3$ is preferred over $\kappa=1$. The reward-temperature response is non-monotone, reflecting the trade-off between broad exploration and concentration on high-quality formulas. The full $36$-configuration grid is reported in \tref{tab:alphaqd_ablation_full}; its purpose is to show the sensitivity of the method rather than to treat every grid cell as a statistically separated comparison.

This grid varies structural hyperparameters rather than removing complete modules. It should therefore be read as a sensitivity analysis, complementary to the archive diagnostics and main comparison, rather than as a component-by-component ablation. The descriptor dimension also changes the number of possible cells, while $K^\star=\min(10,5\kappa)$ links the cell capacity to the exported pool size; these coupled effects are considered when interpreting the marginal trends.

\begin{table}[!htb]
\centering
\caption[AlphaQD hyperparameter sensitivity.]{AlphaQD hyperparameter sensitivity on the A-share validation set under the unified-ridge protocol (single seed). The canonical configuration is $d_b=2$, $\kappa=3$, and $\beta_q=0.5$. Columns are descriptor dimension $d_b$, cell capacity $\kappa$, reward temperature $\beta_q$, signal IC, ICIR, Spearman rank-IC, and net Sharpe.}
\label{tab:alphaqd_ablation_full}
\renewcommand{\arraystretch}{1.05}
\footnotesize
\begin{tabular}{ccccccc}
\toprule
$d_b$ & $\kappa$ & $\beta_q$ & IC & ICIR & RankIC & Sharpe \\
\midrule
$1$ & $1$ & $0.5$ & $0.0565$ & $0.455$ & $0.074$ & $0.322$ \\
$1$ & $1$ & $1.0$ & $0.0584$ & $0.592$ & $0.097$ & $0.951$ \\
$1$ & $1$ & $2.0$ & $0.0772$ & $0.647$ & $0.127$ & $1.027$ \\
$1$ & $1$ & $5.0$ & $0.0328$ & $0.246$ & $0.045$ & $0.889$ \\
$1$ & $3$ & $0.5$ & $0.0740$ & $0.542$ & $0.123$ & $1.186$ \\
$1$ & $3$ & $1.0$ & $0.0577$ & $0.464$ & $0.094$ & $0.699$ \\
$1$ & $3$ & $2.0$ & $0.0757$ & $0.630$ & $0.106$ & $0.540$ \\
$1$ & $3$ & $5.0$ & $0.0760$ & $0.511$ & $0.127$ & $0.835$ \\
$1$ & $5$ & $0.5$ & $0.0669$ & $0.493$ & $0.096$ & $0.833$ \\
$1$ & $5$ & $1.0$ & $0.0658$ & $0.560$ & $0.100$ & $0.879$ \\
$1$ & $5$ & $2.0$ & $0.0683$ & $0.664$ & $0.107$ & $1.010$ \\
$1$ & $5$ & $5.0$ & $0.0486$ & $0.485$ & $0.077$ & $0.668$ \\
\midrule
$2$ & $1$ & $0.5$ & $0.0541$ & $0.461$ & $0.075$ & $0.874$ \\
$2$ & $1$ & $1.0$ & $0.0258$ & $0.218$ & $0.040$ & $0.909$ \\
$2$ & $1$ & $2.0$ & $0.0724$ & $0.662$ & $0.106$ & $1.175$ \\
$2$ & $1$ & $5.0$ & $0.0760$ & $0.646$ & $0.096$ & $0.738$ \\
$2$ & $3$ & $0.5$ & $0.0736$ & $0.537$ & $0.123$ & $0.955$ \\
$2$ & $3$ & $1.0$ & $0.0432$ & $0.366$ & $0.047$ & $0.907$ \\
$2$ & $3$ & $2.0$ & $0.0724$ & $0.541$ & $0.112$ & $0.718$ \\
$2$ & $3$ & $5.0$ & $0.0712$ & $0.565$ & $0.113$ & $0.869$ \\
$2$ & $5$ & $0.5$ & $0.0589$ & $0.529$ & $0.087$ & $0.976$ \\
$2$ & $5$ & $1.0$ & $0.0793$ & $0.613$ & $0.122$ & $0.892$ \\
$2$ & $5$ & $2.0$ & $0.0375$ & $0.511$ & $0.087$ & $1.005$ \\
$2$ & $5$ & $5.0$ & $0.0734$ & $0.578$ & $0.123$ & $1.141$ \\
\midrule
$3$ & $1$ & $0.5$ & $0.0699$ & $0.505$ & $0.109$ & $0.991$ \\
$3$ & $1$ & $1.0$ & $0.0318$ & $0.368$ & $0.058$ & $0.669$ \\
$3$ & $1$ & $2.0$ & $0.0502$ & $0.516$ & $0.080$ & $0.644$ \\
$3$ & $1$ & $5.0$ & $0.0669$ & $0.534$ & $0.083$ & $1.144$ \\
$3$ & $3$ & $0.5$ & $0.0756$ & $0.525$ & $0.127$ & $1.040$ \\
$3$ & $3$ & $1.0$ & $0.0603$ & $0.518$ & $0.078$ & $0.581$ \\
$3$ & $3$ & $2.0$ & $0.0554$ & $0.616$ & $0.085$ & $0.957$ \\
$3$ & $3$ & $5.0$ & $0.0355$ & $0.405$ & $0.048$ & $-0.008$ \\
$3$ & $5$ & $0.5$ & $0.0314$ & $0.352$ & $0.031$ & $0.422$ \\
$3$ & $5$ & $1.0$ & $0.0741$ & $0.562$ & $0.130$ & $1.025$ \\
$3$ & $5$ & $2.0$ & $0.0494$ & $0.414$ & $0.087$ & $0.584$ \\
$3$ & $5$ & $5.0$ & $0.0306$ & $0.342$ & $0.044$ & $0.847$ \\
\bottomrule
\end{tabular}
\end{table}

\paragraph{Computational cost.} A full training run of AlphaQD on the A-share panel---$2{,}000$ GFlowNet trajectory-balance episodes with mini-batch size $64$, archive maintenance, and PFS evaluation---completes in approximately five to six hours on a single NVIDIA RTX~4090 GPU. The dominant cost is the per-candidate backtest of the IC series on the training panel. With a fixed grid and per-cell quota, an ordinary cell lookup and insertion is $O(1)$; periodic PCA updates additionally reproject the stored elites. After training, the exported $K^\star=10$ alphas and validation-fitted ridge combiner are fixed.

%----------------------------------------------------------------------------------------
\section{Chapter Summary}
\label{sec:alphaqd_summary}

This chapter presented AlphaQD, a quality-diversity RL method for formulaic alpha mining. Its central choice is structural: the running pool is removed from the generator's reward, and diversity, robustness, and selection are handled by a MAP-Elites archive whose behaviour descriptor is the normalised daily IC series---an empirical proxy for a formula's return source. The training-window reward is thereby invariant to archive updates. Conditional on a realised archive, the final cell-selection class scales with the number of occupied cells rather than the cardinality of the formula space. On the A-share panel AlphaQD attains the highest annualised return and net Sharpe among learned formula-mining methods; among the four generative/RL pool-style methods it has the lowest turnover on both panels. On the U.S.\ panel it ranks third on IC, ICIR, and annualised return and second on net Sharpe and turnover. In the arc of the thesis, AlphaQD removes the aggregation mechanism from the learning reward entirely, complementing the learned router of EarnHFT and the VAE-based router of FineFT.
